% Esquisse d'un programme (cote 119, copy 1): archive cote 119, site page 45 (PDF page 46), typescript page 43.
% typescript. Draft compared: EsquisseFr.pdf (text layer):1549-1600.
% AI-transcribed by Claude Opus 5.5 (claude-opus-5-5, adaptive thinking, effort high), Message Batches API, batch esq1 (msgbatch_01LNsogFyF3y6KSp9UbtJtBk), 2026-11-12; not yet checked by eye.
\scanpage{119/45}
\mspage{43}

volumes à paraître chez Hermann, sous le titre commun ``\uline{A la Poursuite}
\uline{des Champs}''. Je prévois actuellement (après des élargissements
successifs du propos initial) que le manuscript de l'ensemble des
deux volumes, que j'espère achever en cours d'année pour ne plus
avoir à y revenir, fera dans les 1 500 pages dactylographiées.
Ces deux volumes d'ailleurs sont pour moi les premiers dans une
série plus vaste, sous le titre commun ``\uline{Réflexions Mathématiques}'',
où je compte développer tant soit peu certains des thèmes esquis-%
sés dans le présent rapport.

Vu qu'il s'agit d'un travail en cours de rédaction,
et même d'achèvement, dont le premier volume sans doute paraîtra
cette année et contiendra une introduction circonstanciée, il
est sans doute moins intéressant que je m'étende ici sur ce thème
de réflexion, et me contenterai donc d'en parler très brièvement.
Ce travail me semble quelque peu marginal par rapport aux thèmes
que je viens d'esquisser, et ne représente pas (il me semble) un
véritable renouvellement d'optique ou d'approche par rapport à
mes intérêts et ma vision mathématiques d'avant 1970. Si je m'y
suis résolu soudain, c'est presque en déserpoir de cause, alors
que près de vingt ans se sont écoulés depuis que se sont posées
en termes bien clairs un certain nombre de questions visiblement
fondamentales, et mûres pour être menées à leur terme, sans que
personne ne les voie, ou prenne la peine de les sonder. Aujour-%
d'hui encore, les structures de base qui interviennent dans le
point de vue homotopique en topologie, y compris même en algèbre
homologique commutative, ne sont pas comprises, et à ma connais-%
sance, après les travaux de Verdier, de Giraud et d'Illusie sur
ce thème (qui constituent autant de ``coups d'envoi'' attendant tou-%
jours une suite ...) il n'y a pas eu d'effort dans ce sens. Je
devrais faire exception sans doute pour le travail d'axiomatisa-%
tion fait par Quillen sur la notion de catégorie de modèles, à
la fin des années 60, et repris sous des variantes diverses par
divers auteurs. Ce travail à l'époque, et maintenant encore, m'a
beaucoup séduit et appris, tout en allant dans une direction assez

\margin{45}
